\chapter{Quantum circuits}

\abstract{
    Quantum mechanics is formulated in tensor spaces and quantum circuits are naturally tensor networks.
    Quantum Circuits are naturally tensor networks.
    This motivates the investigation, which contractions can be efficiently performed on a Quantum Computer.
}

\subimport{foundations/}{foundations}
\subimport{directed-circuits/}{directed-circuits}
\subimport{computation-circuits/}{computation-circuits}
\subimport{algorithms/}{algorithms.tex}

\section{Conclusion}

\todo{Update.}

We have seen that preparing some classically hard to compute contraction can be done efficiently by a quantum circuit.
We showed in particular that contracted basis encodings can be prepared by circuits, which are classically hard to contract in the presence of undirected loops in the decomposition hypergraph.
However, when assessing the prepared contraction state by computational basis measurements, we effectively draw samples from the corresponding distribution, which can also be achieved classically.
A quantum advantage could arise from different measurement schemes studied in the quantum state tomography community (see \cite{gross_quantum_2010,cramer_efficient_2010,roth_semi-device-dependent_2023}).
To this end, future research could show, how hard to achieve expectation queries can efficiently be prepared by customized measurement schemes.
